
\begin{proposition}
\label{XIII.4.1}
\setcounter{subsection}{\value{proposition}}% sincroniza la etiqueta visible XIII.4.1.1
Sean $S$ un esquema conexo y $f\colon X\to S$ un morfismo localmente
$0$-acíclico \textup{(SGA~4 XV~1.11)}; supongamos además que~$f$ es
$0$-acíclico (lo que, cuando~$f$ es coherente, equivale a decir que las
fibras geométricas de~$f$ son conexas \textup{(SGA~4 XV~1.16)}). Sea~$\LL$
un conjunto de números primos. Si~$S'$ es un esquema étale sobre~$S$, se
denotan por~$X'$ y~$f'$ las imágenes inversas de~$X$ y~$f$ sobre~$S'$.
Supongamos que, para todo recubrimiento étale~$S'$ de~$S$ y para todo
recubrimiento étale~$E$ de~$X'$, cociente de un recubrimiento galoisiano
cuyo grupo es un~$\LL$-grupo, $(E,f')$ es cohomológicamente propio
relativamente a~$S'$ en dimensión $\leqslant 0$ y~$f'_*E$ es constructible.
Entonces, si~$\bar{s}$ es un punto geométrico de~$S$ y~$a$ un punto
geométrico de la fibra~$X_{\bar{s}}$, la sucesión de homomorfismos de grupos
\begin{equation*}
\label{eq:XIII.4.1.1}
\tag{\thesubsection.1}
\pi_1^{\LL}(X_{\bar{s}},a)\lto{u} \pi_1'(X,a)\lto{v}\pi_1(S,a)\to1
\end{equation*}
es exacta.
\end{proposition}

Este enunciado generaliza X.\Ref{X.1.4}, cuya demostración vamos a repetir.

Demostremos primero que~$v$ es sobreyectivo. Basta demostrar que, para todo
recubrimiento étale conexo~$S'$ de~$S$, $X'$ también es conexo
(\Ref{V.6.9}). Sea~$C$ un conjunto con al menos dos elementos. Del hecho
de que~$f$ es $0$-acíclico resulta que el morfismo canónico
$$
\H^0(S',C_{S'})\to \H^0(X',C_{X'})
$$
es biyectivo; por tanto, $X'$ es conexo, de donde se sigue la sobreyectividad
de~$v$.

Por
\marginpar{416}
definición de~$K^{\LL}$ \eqref{XIII.4.0}, se tiene la sucesión exacta
$$
1\to K^{\LL} \to \pi_1'(X,a)\to \pi_1(S,a)\to 1
$$
Sean $\widetilde{S}$ el recubrimiento universal de~$S$ y
$\widetilde{X}=\widetilde{S}\times_SX$; el grupo~$K^{\LL}$ clasifica los
recubrimientos galoisianos~$P$ cuyo grupo es un~$\LL$-grupo y para los
cuales existen un recubrimiento étale~$S'$ de~$S$ y un recubrimiento
galoisiano~$Q$ de~$X'=X\times_SS'$ tales que se tiene un isomorfismo
$P\simeq Q\times_{X'}\widetilde{X}$. Para que la sucesión
\eqref{eq:XIII.4.1.1} sea exacta, es necesario y suficiente que el morfismo
canónico
$$
\pi_1^{\LL}(X_{\bar{s}},a)\to K^{\LL}
$$
sea sobreyectivo. Según la interpretación de~$K^{\LL}$, esto equivale a
decir que, para todo recubrimiento étale~$S'$ de~$S$ y para todo
recubrimiento galoisiano~$Q$ de~$X'$ cuyo grupo es un~$\LL$-grupo, si
$P=Q\times_{X'}\widetilde{X}$ es conexo, entonces $Q|X_{\bar{s}}$ es
conexo. Demostremos que se cumple esta última condición. Sean, en efecto,
$S'$ un recubrimiento étale de~$S$ y~$Q$ un recubrimiento galoisiano de~$X'$
cuyo grupo es un~$\LL$-grupo~$F$, tal que~$Q|X_{\bar{s}}$ es disconexo, y
demostremos que, tras reemplazar~$S'$ por un recubrimiento étale, $Q$ pasa
a ser disconexo. Existe un subgrupo~$G$ de~$F$, distinto de~$F$, y un
torsor~$R$ bajo~$G|X_{\bar{s}}$ tal que~$Q|X_{\bar{s}}$ se obtiene, por
extensión del grupo estructural $G\to F$, a partir de~$R$. El recubrimiento
étale
$E=Q/G$ de~$X'$ es tal que~$E|X_{\bar{s}}$ tiene una sección. Según
\Ref{XIII.1.16},
\ifthenelse{\boolean{orig}}
{$f_*E$}
{$f'_*E$}
es localmente constante y constructible y, tras reemplazar~$S'$ por un
recubrimiento étale, puede incluso suponerse que
\ifthenelse{\boolean{orig}}
{$f_*E$}
{$f'_*E$}
es constante.
Como~$(E,f')$ es cohomológicamente propio relativamente a~$S'$ en dimensión
$\leqslant0$ y~$\H^0(X_{\bar{s}},E|X_{\bar{s}})$ no es vacío, se ve que~$E$
tiene una sección. Pero esto demuestra que~$Q$ es disconexo, lo que
concluye la demostración.

De~\Ref{XIII.4.1} se deduce el lema siguiente, que se utilizará
en~\Ref{XIII.4.6}.
